\subsection*{Question 1 \quad Réponse D}

On applique l'identité remarquable $(a-b)^2 = a^2 - 2ab + b^2$ avec $a = 5x$ et $b = 2$ :
\[(5x-2)^2 = (5x)^2 - 2\times 5x\times 2 + 2^2 = 25x^2 - 20x + 4.\]

\subsection*{Question 2 \quad Réponse D}

La multiplication est prioritaire sur la soustraction :
\[\dfrac{1}{4}\times\dfrac{8}{5} - \dfrac{1}{5} = \dfrac{8}{20} - \dfrac{1}{5} = \dfrac{2}{5} - \dfrac{1}{5} = \dfrac{1}{5}.\]

\subsection*{Question 3 \quad Réponse A}

Multiplier par $\dfrac{3}{4} = 0,75$ revient à appliquer le coefficient multiplicateur
$0,75 = 1 - 0,25$ : c'est une baisse de 25\,\%.

\subsection*{Question 4 \quad Réponse B}

Un produit de facteurs est nul si et seulement si l'un au moins de ses facteurs est nul :
\[(2x-4)(-3x+3) = 0 \iff 2x - 4 = 0 \ \text{ ou } \ -3x + 3 = 0 \iff x = 2 \ \text{ ou } \ x = 1.\]
D'où $S = \{2 ; 1\}$.

\subsection*{Question 5 \quad Réponse C}

Une baisse de 50\,\% correspond au coefficient $0,5$ et une hausse de 50\,\% au
coefficient $1,5$. Le coefficient multiplicateur global est leur produit :
\[0,5 \times 1,5 = 0,75 = 1 - 0,25.\]
Le prix a donc diminué de 25\,\%.

\subsection*{Question 6 \quad Réponse B}

Le coefficient directeur de la droite (AB) vaut
\[a = \dfrac{y_B - y_A}{x_B - x_A} = \dfrac{3 - (-1)}{1 - 3} = \dfrac{4}{-2} = -2.\]
On a donc $f(x) = -2x + b$, et comme B$(1~;~3)$ appartient à la droite :
$3 = -2\times 1 + b$, soit $b = 5$. D'où $f(x) = -2x + 5$.

\subsection*{Question 7 \quad Réponse C}

La moyenne coefficientée est la somme des notes multipliées par leur coefficient, divisée par
la somme des coefficients :
\[\overline{x} = \dfrac{18\times 1 + 11\times 2 + 16\times 1}{1 + 2 + 1} = \dfrac{18 + 22 + 16}{4} = \dfrac{56}{4} = 14.\]

\subsection*{Question 8 \quad Réponse D}

6,4~h, c'est 6~h plus 0,4~h. Or une heure compte 60 minutes :
\[0,4 \times 60 = 24.\]
Donc 6,4~h $=$ 6~h 24~min.
